\documentclass[10pt]{article}

\usepackage[margin=0.72in]{geometry}
\usepackage{amsmath,amssymb}
\usepackage{booktabs}
\usepackage{microtype}

\setlength{\parindent}{0pt}
\setlength{\parskip}{0.45em}
\newcommand{\E}{\mathcal{F}}

\title{Why Normalization Does Not Force $s_1$ and $s_2$ to Unity}
\author{}
\date{}

\begin{document}
\maketitle
\vspace{-2.2em}

\paragraph{Question.}
The rectangular forward initializes two per-experiment coefficients at one and
uses fixed probe and detector scaling. Does that construction force
$s_1,s_2\approx1$, leaving them only a small modulation such as $0.9$--$1.1$?

\paragraph{Short answer.}
No. The fixed factors cancel known preprocessing and choose detector/probe
units; they do not normalize the unknown object coefficients. Unity is an
initialization (or a plausible scale under a well-matched gauge), not an
invariant or bound.

\begin{center}
\small
\begin{tabular}{@{}ll@{}}
\toprule
Symbol & Meaning \\
\midrule
$P_{\rm phys}$ & complex probe in the chosen acquisition gauge \\
$q$ & fixed probe normalization, $P_{\rm train}=qP_{\rm phys}$ \\
$d$ & fixed legacy loader intensity-normalization metadata \\
$k$ & derived exit-wave amplitude factor passed to the forward as \texttt{scale} \\
$a,b$ & learned dimensionless real/imaginary decoder textures \\
$s_1,s_2$ & per-experiment real coefficients multiplying those textures \\
\bottomrule
\end{tabular}
\end{center}

\section*{What the fixed scales actually cancel}

The autograd rectangular forward computes
\begin{equation}
  \psi = k\left[s_1(P_{\rm train}a)
           + i\,s_2(P_{\rm train}b)\right],
  \qquad
  I_{\rm pred}=\sum_p\left|\E\{\psi_p\}\right|^2 .
\end{equation}
Here $k$ is a real scalar; it is not the complex probe, which is the separate
factor $P_{\rm train}$.

In the loader-normalized legacy implementation, the notation maps directly to
\begin{align*}
  q &\longleftrightarrow \texttt{batch[2]},\\
  d &\longleftrightarrow
      \texttt{fields['physics\_scaling\_constant']}.
\end{align*}
The loss code constructs
\begin{equation}
  \texttt{exit\_wave\_scale}
  \equiv k=(q^2d+10^{-9})^{-1/2}
\end{equation}
under the variable name \texttt{output\_scale}; the rectangular forward then
receives that value through its argument named \texttt{scale}.

For the current \texttt{ci\_intensity\_v2} count-intensity path,
\begin{equation}
  k=q^{-1}, \qquad kP_{\rm train}=q^{-1}(qP_{\rm phys})=P_{\rm phys}.
\end{equation}
Thus the output factor exactly undoes training-probe normalization. There is no
\texttt{physics\_scale} in this CI path; measured intensities remain counts.

For the explicitly loader-normalized \texttt{legacy\_v1} rectangular path,
\begin{equation}
  k=(q^2d+10^{-9})^{-1/2}
  \approx \frac{1}{q\sqrt d},
  \qquad kP_{\rm train}\approx\frac{P_{\rm phys}}{\sqrt d}.
\end{equation}
This statement is specific to that adapter; other legacy adapters, including
dictionary parity, need not supply the same factors. In either path, the fixed
terms remove known representation choices. They leave the unknown object
\begin{equation}
  O=s_1a+i\,s_2b
\end{equation}
inside the physical forward. Nothing in the cancellation sets either
coefficient to one.

\section*{Why the coefficients are not identified near one}

First, $s_1$ and $s_2$ describe real and imaginary object contrast, not merely
residual calibration. Those contrasts need not be equal or have unit scale.

Second, the model has an internal factorization freedom. Wherever the rescaled
texture remains in the decoder's representable range,
\begin{equation}
  s_1a=(cs_1)(a/c), \qquad
  s_2b=(rs_2)(b/r).
\end{equation}
The detector loss observes the products, not each factor separately. A bounded
decoder limits this freedom but does not fix a texture norm, so it does not
select $s_1=s_2=1$.

Third, ptychography has a distinct upstream acquisition gauge:
\begin{equation}
  P_{\rm phys}O=(\gamma P_{\rm phys})(O/\gamma).
\end{equation}
Storing a physical probe chooses one such decomposition; diffraction alone
does not establish absolute object units. This acquisition gauge is separate
from the train-time decoder factorization above.

\section*{Initialization versus construction}

The module constructor creates $s_1=s_2=1$. With \texttt{ones}
initialization, training starts there. With \texttt{dose\_closure}, both are
instead set to a data-derived common gauge $g>0$ so a unit-object prediction
starts near the measured dose; $g$ need not equal one. During autograd training
the coefficients are ordinary unconstrained parameters: there is no range,
clamp, or regularizer keeping them near unity.

Correct normalization can therefore make order-unity values plausible and
improve conditioning. It cannot guarantee them. Values far from one may
reflect legitimate object contrast, decoder rescaling, the chosen
probe--object gauge, or inconsistent scaling metadata.

\paragraph{Assessment of the quoted interpretation.}
The derived output factor is indeed passed as the forward argument named
\texttt{scale}, and fixed factors do absorb known preprocessing. However,
that \texttt{scale} is not the complex probe, and ``$s_1,s_2$ stay around one
by construction'' is false. At most, near-one behavior is an empirical
expectation under a compatible gauge and training trajectory.

\end{document}
